\documentclass[11pt,a4paper]{article}
\usepackage[T1]{fontenc}
\usepackage[utf8]{inputenc}
\usepackage{amsmath,amssymb,amsthm}
\usepackage[margin=1in]{geometry}
\usepackage[colorlinks=true,linkcolor=blue,urlcolor=blue]{hyperref}
\theoremstyle{plain}
\newtheorem{theorem}{Theorem}
\newtheorem{lemma}[theorem]{Lemma}
\newtheorem{proposition}[theorem]{Proposition}
\newtheorem{corollary}[theorem]{Corollary}
\theoremstyle{definition}
\newtheorem{remark}[theorem]{Remark}

\title{The empirical closed form for OEIS A224207, proved}
\author{Adrian Perez Fontelles\\ \small Independent researcher}
\date{1 September 2026}
\begin{document}
\maketitle

\begin{abstract}
OEIS A224207 carries an empirical closed form for $a(n)$ --- not a recurrence relating
consecutive terms but an explicit formula. It is true, and it is decidable rather than
empirical. The entry counts arrays subject to a condition on a bounded window of consecutive lines, so the lines of an array are the
vertices of a finite digraph, an array is a walk in it, and $a(n)$ is a walk count on
$S=32768$ vertices. Any function of the form $\sum_j p_j(n)\lambda_j^{\,n}$ satisfies the
monic linear recurrence whose characteristic polynomial is
$\prod_j(x-\lambda_j)^{\deg p_j+1}$; here that polynomial is $q(x)=(x-1)^{31}$, of degree
$31$. Two things are then checked exactly: that the walk count satisfies the same
recurrence from some index on, which the Cayley--Hamilton theorem reduces to a finite
computation, and that the two sequences agree at $31$ consecutive indices past that
point. A monic recurrence determines every later term from its predecessors, so agreement
there is agreement for good.
\end{abstract}

\noindent\small 2020 Mathematics Subject Classification. 05A15, 11B37, 15A18.\normalsize

\section{The conjecture}

OEIS A224207 is ``Number of 5Xn 0..3 arrays with rows unimodal and antidiagonals nondecreasing.''. It has offset $1$ and begins
\[
1024, 160000, 4357084, 54177872, 451319098, 2981774784, 17178450136, 91585310015,\ \dots
\]
The entry states, as a conjecture contributed by its author and never marked settled:
\begin{quote}
Empirical polynomial of degree 30 (see link above)

Empirical: a(n) = (1/7353420799762759680000000)*n\^{}30 + (13/490228053317517312000000)*n\^{}29 + (1787413/659931481843536494592000000)*n\^{}28 + (523777/2856846241746911232000000)*n\^{}27 + (5209787761/576130658752293765120000000)*n\^{}26 + (7095698399/20681613391107981312000000)*n\^{}25 + (16822667231/1618561047999755059200000)*n\^{}24 + (3205378345289/12408968034664788787200000)*n\^{}23 + (6214555829939/1156115034285539328000000)*n\^{}22 + (4667951660587/49047304484841062400000)*n\^{}21 + (925683935693/636977980322611200000)*n\^{}20 + (28808763911323/1486281954086092800000)*n\^{}19 + (8864767135811089/38721556172242944000000)*n\^{}18 + (483369391460041/198572082934579200000)*n\^{}17 + (1551337008095993947/65143323913302835200000)*n\^{}16 + (4793120246036954171/21714441304434278400000)*n\^{}15 + (16473765719932522583/8239854959271936000000)*n\^{}14 + (4257346719396442537/236631732163706880000)*n\^{}13 + (96182338096152968435179/606960392999908147200000)*n\^{}12 + (5979235820735826764833/4598184795453849600000)*n\^{}11 + (288084831281705855693141/31351259969003520000000)*n\^{}10 + (26446074543700186965437/522520999483392000000)*n\^{}9 + (36578999268001533954677/190761634732032000000)*n\^{}8 + (2929751725546871424717823/7344322937183232000000)*n\^{}7 + (5529704418386921739771359/22440986752504320000000)*n\^{}6 + (2930125988920364980531/5394467969352000000)*n\^{}5 + (7349189512583334584671/1944885518550374400)*n\^{}4 - (272138048304685248839/27012298868755200)*n\^{}3 - (4783077062534268221/506830141818000)*n\^{}2 + (1232877264268367/21565644100)*n - 54467 for n>3
\end{quote}
As of the ``Last modified'' line on the live entry (Oct 03 2025 23:19:46, revision 10) this is still
recorded as empirical, and nothing on the entry records it as proved. Write
\[
f(n)\;=\;\frac{n^{30}}{7353420799762759680000000} + \frac{13 n^{29}}{490228053317517312000000} + \frac{1787413 n^{28}}{659931481843536494592000000} + \frac{523777 n^{27}}{2856846241746911232000000} + \frac{5209787761 n^{26}}{576130658752293765120000000} + \frac{7095698399 n^{25}}{20681613391107981312000000} + \frac{16822667231 n^{24}}{1618561047999755059200000} + \frac{3205378345289 n^{23}}{12408968034664788787200000} + \frac{6214555829939 n^{22}}{1156115034285539328000000} + \frac{4667951660587 n^{21}}{49047304484841062400000} + \frac{925683935693 n^{20}}{636977980322611200000} + \frac{28808763911323 n^{19}}{1486281954086092800000} + \frac{8864767135811089 n^{18}}{38721556172242944000000} + \frac{483369391460041 n^{17}}{198572082934579200000} + \frac{1551337008095993947 n^{16}}{65143323913302835200000} + \frac{4793120246036954171 n^{15}}{21714441304434278400000} + \frac{16473765719932522583 n^{14}}{8239854959271936000000} + \frac{4257346719396442537 n^{13}}{236631732163706880000} + \frac{96182338096152968435179 n^{12}}{606960392999908147200000} + \frac{5979235820735826764833 n^{11}}{4598184795453849600000} + \frac{288084831281705855693141 n^{10}}{31351259969003520000000} + \frac{26446074543700186965437 n^{9}}{522520999483392000000} + \frac{36578999268001533954677 n^{8}}{190761634732032000000} + \frac{2929751725546871424717823 n^{7}}{7344322937183232000000} + \frac{5529704418386921739771359 n^{6}}{22440986752504320000000} + \frac{2930125988920364980531 n^{5}}{5394467969352000000} + \frac{7349189512583334584671 n^{4}}{1944885518550374400} - \frac{272138048304685248839 n^{3}}{27012298868755200} - \frac{4783077062534268221 n^{2}}{506830141818000} + \frac{1232877264268367 n}{21565644100} - 54467.
\]

\section{The count is a walk count}

The entry's defining condition is local: it is a condition on a bounded window of consecutive lines. Take as vertices the
admissible configurations of such a window and put an edge from $u$ to $v$ when $v$ may
follow $u$, which the condition decides by inspection. An admissible array is then exactly a
walk, so with $M$ the adjacency matrix of that digraph on $S=32768$ vertices, and $u,w$ the
vectors recording which windows may start and end an array,
\[
a(n)\;=\;u^{\!\top}M^{\,g(n)}w
\]
for the linear function $g$ that the entry's shape supplies. In particular $a$ is
$C$-finite: it satisfies some monic integer linear recurrence, though not necessarily a short
one.

\section{A closed form is a recurrence}

\begin{lemma}\label{lem:ann}
Let $f(m)=\sum_{j}p_j(m)\lambda_j^{\,m}$ with the $\lambda_j$ distinct and the $p_j$
polynomials, and put $q(x)=\prod_j(x-\lambda_j)^{\deg p_j+1}$. Then $q$ applied to the shift
operator annihilates $f$: writing $q(x)=x^{r}-\sum_{i=1}^{r}c_ix^{r-i}$,
\[
f(m)\;=\;\sum_{i=1}^{r}c_i\,f(m-i)\qquad\text{for every }m.
\]
\end{lemma}

\begin{proof}
The shift operator $E$ acts on $m\mapsto\lambda^{m}p(m)$ as
$(E-\lambda)\bigl(\lambda^{m}p(m)\bigr)=\lambda^{m+1}\bigl(p(m+1)-p(m)\bigr)$, and
$p(m+1)-p(m)$ has degree one less than $p$. So $(E-\lambda)^{\deg p+1}$ kills
$\lambda^{m}p(m)$, and the factors of $q$ commute.
\end{proof}

Here $q(x)=(x-1)^{31}$, of degree $31$; the closed form is a polynomial in $n$, so this is $(x-1)^{31}$, and the recurrence it encodes is
\begin{equation}\label{eq:rec}
a(n)\;=\;31\,a(n-1) - 465\,a(n-2) + 4495\,a(n-3) - 31465\,a(n-4) + 169911\,a(n-5) - 736281\,a(n-6) + 2629575\,a(n-7) - 7888725\,a(n-8) + 20160075\,a(n-9) - 44352165\,a(n-10) + 84672315\,a(n-11) - 141120525\,a(n-12) + 206253075\,a(n-13) - 265182525\,a(n-14) + 300540195\,a(n-15) - 300540195\,a(n-16) + 265182525\,a(n-17) - 206253075\,a(n-18) + 141120525\,a(n-19) - 84672315\,a(n-20) + 44352165\,a(n-21) - 20160075\,a(n-22) + 7888725\,a(n-23) - 2629575\,a(n-24) + 736281\,a(n-25) - 169911\,a(n-26) + 31465\,a(n-27) - 4495\,a(n-28) + 465\,a(n-29) - 31\,a(n-30) + a(n-31).
\end{equation}

\section{The proof}

\begin{lemma}\label{lem:crit}
With $q$ as above, $a$ satisfies \eqref{eq:rec} for all large $n$ if and only if
$u^{\!\top}M^{\,j}q(M)w=0$ for every $j\ge0$, and it is enough to check $j<S$.
\end{lemma}

\begin{proof}
Substituting the walk expression, the residual $a(n)-\sum_ic_ia(n-i)$ equals
$u^{\!\top}M^{\,g(n)-r}q(M)w$ up to the fixed linear reparametrisation $g$. For the bound,
$M^{S}$ is an integer combination of $I,M,\dots,M^{S-1}$ by Cayley--Hamilton, so the values
$u^{\!\top}M^{j}q(M)w$ satisfy the monic recurrence given by the characteristic polynomial of
$M$; $S$ consecutive zeros therefore force all later ones, and the last nonzero value pins the
threshold exactly.
\end{proof}

\begin{theorem}
$a(n)=f(n)$ for every $n\ge4$.
\end{theorem}

\begin{proof}
The vector $w'=q(M)w$ was formed by $31$ matrix--vector products in exact integer
arithmetic, and $u^{\!\top}M^{j}w'$ evaluated for $j=0,1,\dots$, again exactly. Those integers
vanish from the point corresponding to $n=35$ onwards, and the run was continued until the
working vector was identically zero, which settles every larger $j$ at once. So $a$ satisfies
\eqref{eq:rec} for every $n>34$. By Lemma~\ref{lem:ann} $f$ satisfies \eqref{eq:rec}
identically. Both were then evaluated in exact arithmetic at the $31$ consecutive indices
$n=35,\dots,65$ and agree at each. A monic recurrence of order $31$
determines $a(m)$ from $a(m-1),\dots,a(m-31)$, and for every $m\ge66$ both
$m>34$ and $m-31\ge35$ hold, so induction on $m$ gives $a(m)=f(m)$ for every
$m\ge35$.
Below $n=35$ the recurrence argument gives nothing, since \eqref{eq:rec} is only
established for $a$ from $n=35$ on. The remaining indices $n=4,\dots,34$ were
therefore checked one at a time, directly against the walk count in exact integer arithmetic,
and agree.
\end{proof}

The range is tight in the sense that was checked: at $n=3$ the closed form and the sequence differ, so no larger range is available.

\section{Verification}

Four checks, each independent of the algebra above.

First, the walk model reproduces every one of the $16$ terms the entry publishes,
exactly, in integer arithmetic. This is what ties the model to the entry: the model is built
by reading the entry's English, and a misreading gives a different digraph and different
counts.

Second, the closed form was evaluated in exact rational arithmetic --- not floating point ---
at every published term from $n=4$ on, and agrees with the entry's own DATA at each.

Third, the annihilation test of Lemma~\ref{lem:crit} was carried out exactly, and the model
and the closed form were then compared directly at sixty consecutive indices beyond
$n=4$, far past where the proof already settles the matter.

Fourth, the closed form was deliberately perturbed --- by a constant and by a term linear in
$n$ --- and each perturbed form was rejected by the same comparison, so the test is not one
that anything would pass.

\begin{thebibliography}{9}
\bibitem{oeis} The OEIS Foundation, \emph{The On-Line Encyclopedia of Integer
Sequences}, \url{https://oeis.org}, 2026. Sequence A224207.
\bibitem{stanley} R.~P.~Stanley, \emph{Enumerative Combinatorics}, Volume 1, 2nd ed.,
Cambridge University Press, 2011. (Transfer-matrix method, Section 4.7; $C$-finite sequences,
Section 4.1.)
\bibitem{kp} M.~Kauers and P.~Paule, \emph{The Concrete Tetrahedron}, Springer, 2011.
(Closed forms and linear recurrences with constant coefficients.)
\bibitem{horn} R.~A.~Horn and C.~R.~Johnson, \emph{Matrix Analysis}, 2nd ed., Cambridge
University Press, 2013.
\end{thebibliography}

\end{document}
